% --- 02 HQC KEM in one slide -------------------------------------------------
\begin{frame}{HQC: ανταλλαγή κλειδιού βασισμένη σε κώδικες διόρθωσης σφαλμάτων}
  \begin{columns}[T,onlytextwidth]
    \begin{column}{0.44\textwidth}
      \textbf{Τι είναι}
      \begin{itemize}
        \item \en{Code-based} \KEM στον δακτύλιο $\GFquot{N}$
        \item Συνενωμένος κώδικας: \en{Reed--Solomon} (εξωτ.) $\circ$ \en{Reed--Muller} (εσωτ.)
        \item 3 επίπεδα ασφαλείας: \HQC-128/192/256 (\en{NIST} 1/3/5)
      \end{itemize}
      \vspace{0.4em}
      \textbf{Τρεις κυρίαρχες πράξεις}
      \begin{itemize}
        \item \HL{polymul} στον δακτύλιο
        \item \HL{Keccak} (\en{SHA-3}/\en{SHAKE}): hashing \& seed expansion
        \item \HL{decode} (\en{RM} + \en{RS})
      \end{itemize}
    \end{column}
    \begin{column}{0.56\textwidth}
      \resizebox{\linewidth}{!}{% HQC KEM sequence diagram: Alice ↔ Bob key exchange
% Ops moved fully outside the lifelines + activation bars (U1) + primitive tags (U2)
% Usage: % HQC KEM sequence diagram: Alice ↔ Bob key exchange
% Usage: % HQC KEM sequence diagram: Alice ↔ Bob key exchange
% Usage: \input{figures/hqc-kem-sequence}
\begin{tikzpicture}[
  entity/.style={draw, rounded corners=3pt, fill=thesisblue!8,
                 minimum width=2.2cm, minimum height=0.8cm, font=\sffamily\bfseries},
  msg/.style={->, >=stealth, thick},
  op/.style={font=\small, align=left},
  note/.style={font=\scriptsize\itshape, text=gray!70!black},
]

% Entities
\node[entity] (alice) at (0, 0) {Alice};
\node[entity] (bob) at (8, 0) {Bob};

% Lifelines
\draw[dashed, gray!50] (alice.south) -- ++(0, -9.5);
\draw[dashed, gray!50] (bob.south) -- ++(0, -9.5);

% KeyGen (Bob)
\node[op, right] at (6.5, -1.2) {\KeyGen};
\node[note, right] at (6.5, -1.7) {$(ek, dk) \leftarrow \KeyGen()$};

% Send public key
\draw[msg, <-] (0, -2.5) -- node[above, font=\small] {$ek$} (8, -2.5);

% Encaps (Alice)
\node[op, left] at (1.5, -3.5) {\Encaps};
\node[note, left] at (1.5, -4.0) {$(ct, ss) \leftarrow \Encaps(ek)$};

% Send ciphertext
\draw[msg] (0, -5.0) -- node[above, font=\small] {$ct$} (8, -5.0);

% Decaps (Bob)
\node[op, right] at (6.5, -6.0) {\Decaps};
\node[note, right] at (6.5, -6.5) {$ss' \leftarrow \Decaps(dk, ct)$};

% Shared secret
\node[font=\small\bfseries, text=green!50!black] at (0, -7.8) {$ss$};
\node[font=\small\bfseries, text=green!50!black] at (8, -7.8) {$ss' = ss$};
\draw[thick, green!50!black, dashed] (0.5, -7.8) -- (7.5, -7.8);
\node[note, above] at (4, -7.6) {κοινό μυστικό};

% Annotations: which kernel is called where
\node[note, left, text=red!60!black] at (-0.3, -4.5) {polymul, Keccak};
\node[note, right, text=red!60!black] at (8.3, -6.8) {polymul, Keccak, decode};

\end{tikzpicture}

\begin{tikzpicture}[
  entity/.style={draw, rounded corners=3pt, fill=thesisblue!8,
                 minimum width=2.2cm, minimum height=0.8cm, font=\sffamily\bfseries},
  msg/.style={->, >=stealth, thick},
  op/.style={font=\small, align=left},
  note/.style={font=\scriptsize\itshape, text=gray!70!black},
]

% Entities
\node[entity] (alice) at (0, 0) {Alice};
\node[entity] (bob) at (8, 0) {Bob};

% Lifelines
\draw[dashed, gray!50] (alice.south) -- ++(0, -9.5);
\draw[dashed, gray!50] (bob.south) -- ++(0, -9.5);

% KeyGen (Bob)
\node[op, right] at (6.5, -1.2) {\KeyGen};
\node[note, right] at (6.5, -1.7) {$(ek, dk) \leftarrow \KeyGen()$};

% Send public key
\draw[msg, <-] (0, -2.5) -- node[above, font=\small] {$ek$} (8, -2.5);

% Encaps (Alice)
\node[op, left] at (1.5, -3.5) {\Encaps};
\node[note, left] at (1.5, -4.0) {$(ct, ss) \leftarrow \Encaps(ek)$};

% Send ciphertext
\draw[msg] (0, -5.0) -- node[above, font=\small] {$ct$} (8, -5.0);

% Decaps (Bob)
\node[op, right] at (6.5, -6.0) {\Decaps};
\node[note, right] at (6.5, -6.5) {$ss' \leftarrow \Decaps(dk, ct)$};

% Shared secret
\node[font=\small\bfseries, text=green!50!black] at (0, -7.8) {$ss$};
\node[font=\small\bfseries, text=green!50!black] at (8, -7.8) {$ss' = ss$};
\draw[thick, green!50!black, dashed] (0.5, -7.8) -- (7.5, -7.8);
\node[note, above] at (4, -7.6) {κοινό μυστικό};

% Annotations: which kernel is called where
\node[note, left, text=red!60!black] at (-0.3, -4.5) {polymul, Keccak};
\node[note, right, text=red!60!black] at (8.3, -6.8) {polymul, Keccak, decode};

\end{tikzpicture}

\tikzsetnextfilename{ch2-background-hqc-kem-sequence}
\begin{tikzpicture}[
  entity/.style={draw=thesisblue!70, rounded corners=3pt, fill=thesisblue!10,
                 minimum width=2.3cm, minimum height=0.85cm, font=\sffamily\bfseries},
  life/.style={dashed, gray!50},
  msg/.style={-{Stealth[length=2.5mm]}, thick, thesisblue!80},
  act/.style={draw=thesisblue!65, fill=thesisblue!15, rounded corners=1pt},
  op/.style={font=\small\bfseries, text=thesisblue!85!black},
  noteR/.style={font=\scriptsize, align=left},
  noteL/.style={font=\scriptsize, align=right},
  prim/.style={font=\scriptsize\itshape, text=red!65!black},
]

% Actors
\node[entity] (alice) at (0,0) {Alice};
\node[entity] (bob)   at (7,0) {Bob};

% Lifelines
\draw[life] (alice.south) -- ++(0,-8.4);
\draw[life] (bob.south)   -- ++(0,-8.4);

% KeyGen (Bob): activation on the lifeline, label to the RIGHT of Bob
\draw[act] (6.91,-0.75) rectangle (7.09,-1.95);
\node[op,    anchor=west] at (7.35,-1.0)  {\KeyGen};
\node[noteR, anchor=west] at (7.35,-1.45) {$(ek, dk) \leftarrow \KeyGen()$};
\node[prim,  anchor=west] at (7.35,-1.85) {polymul, Keccak};

% Send public key: Bob -> Alice
\draw[msg] (7,-2.4) -- node[above, font=\small] {$ek$} (0,-2.4);

% Encaps (Alice): activation on the lifeline, label to the LEFT of Alice
\draw[act] (-0.09,-2.85) rectangle (0.09,-4.05);
\node[op,    anchor=east] at (-0.35,-3.1)  {\Encaps};
\node[noteL, anchor=east] at (-0.35,-3.55) {$(ct, ss) \leftarrow \Encaps(ek)$};
\node[prim,  anchor=east] at (-0.35,-3.95) {polymul, Keccak};

% Send ciphertext: Alice -> Bob
\draw[msg] (0,-4.6) -- node[above, font=\small] {$ct$} (7,-4.6);

% Decaps (Bob): activation on the lifeline, label to the RIGHT of Bob
\draw[act] (6.91,-5.35) rectangle (7.09,-6.55);
\node[op,    anchor=west] at (7.35,-5.6)  {\Decaps};
\node[noteR, anchor=west] at (7.35,-6.05) {$ss' \leftarrow \Decaps(dk, ct)$};
\node[prim,  anchor=west] at (7.35,-6.45) {polymul, Keccak, decode};

% Shared secret
\node[font=\small\bfseries, text=green!50!black] at (0,-7.6) {$ss$};
\node[font=\small\bfseries, text=green!50!black] at (7,-7.6) {$ss' = ss$};
\draw[dashed, thick, green!50!black] (0.5,-7.6) -- (6.5,-7.6);
\node[font=\scriptsize\itshape, text=gray!70!black, above] at (3.5,-7.5) {κοινό μυστικό};

\end{tikzpicture}
}
    \end{column}
  \end{columns}
  \pdfpcnote{budget ≈75s | exit ≈2:15}
  \pdfpcnote{}
  \pdfpcnote{- KEM → εγκαθιστά κοινό μυστικό → μετά συμμετρική κρυπτο}
  \pdfpcnote{- Bob δημοσιεύει ek → Alice ενθέτει ελεγχόμενο θόρυβο (ct) → Bob αποκωδικοποιεί με dk}
  \pdfpcnote{- polymul, Keccak, decode → οι 3 πράξεις που κυριαρχούν στον χρόνο}
  \pdfpcnote{}
  \pdfpcnote{cut: παρέλειψε το «τι είναι KEM», κράτα μόνο τις 3 πράξεις}
\end{frame}
